\documentclass[letterpaper,10pt]{article}
\usepackage[sinmarca]{utfsm}
\usepackage{tabularx}
\datos{titulo={Cubicación independiente y conciliación del presupuesto tipo}, ramo={Sede Social Tipo C, albañilería armada},
  autor={Nicolás Albornoz Pavez}, fecha={20 de noviembre de 2026}, pie={Ejercicio de portafolio, no oficial}}
% Documento profesional, no universitario: encabezado sin logos institucionales.
\fancyhf{}
\fancyhead[L]{\small Cubicación independiente: Sede Social Tipo C (FSPR 2023)}
\fancyhead[R]{\small Nicolás Albornoz Pavez}
\fancyfoot[L]{\footnotesize Ejercicio de portafolio, no oficial \textperiodcentered\ versión 5}
\fancyfoot[R]{\footnotesize Pág.\ \thepage}
\newcommand{\vLeje}{37.47}
\newcommand{\vDentro}{10}
\newcommand{\vTotal}{20}
\newcommand{\vSob}{376}
\newcommand{\vFal}{423}
\newcommand{\vLad}{88.56}
\newcommand{\vLadPilar}{83.24}
\newcommand{\vFcad}{194.8}
\newcommand{\vTens}{136.3}
\newcommand{\vEsc}{67.1}
\newcommand{\vEscKgm}{0.231}
\newcommand{\vLesc}{290.8}
\newcommand{\vHcad}{1.69}
\newcommand{\vMcad}{23.77}
\newcommand{\vTecho}{89.90}
\newcommand{\vRell}{7.65}

\tikzset{muro/.style={line width=0.55pt}, elem/.style={usmazul,line width=0.2pt}}
\begin{document}

\begin{center}
{\LARGE\bfseries Cubicación independiente y conciliación del presupuesto tipo\par}
\vspace{2pt}{\large Sede Social Tipo C en albañilería armada, Fondo Social Presidente de la República (2023)\par}
\vspace{4pt}{\small Nicolás Albornoz Pavez \textperiodcentered\ Estudiante de Ingeniería Civil, UTFSM \textperiodcentered\ 20 de noviembre de 2026\par}
\end{center}

\section{Resumen}
Se cubicó la obra gruesa y parte de las terminaciones desde el plano vectorial y las especificaciones técnicas (EETT), con los criterios de NCh353.Of2000 (cantidades sin pérdidas; fierro con traslapos, patas y penetración en fundación; vanos de albañilería por tabla 4.a), y se comparó con el presupuesto tipo oficial. \vDentro{} de \vTotal{} partidas quedan dentro de $\pm5$\,\%. Cada una de las restantes se explicó reconstruyendo la cifra del presupuesto: todas, salvo el cielo, se reproducen dentro de $\pm2$\,\% con un criterio identificable.

\begin{resultado}
\begin{itemize}[nosep]
\item \textbf{Criterio de cubicación} (diferencias legítimas): ladrillo sin descuento de vanos menores, sobreexcavación de 5\,\%, moldajes con holgura, tensores todos de $\phi$10 y cubierta sin traslapos.
\item \textbf{Errores del presupuesto}: escalerilla con el doble de la masa del alambre de \SI{4}{mm}, cadena sin estribos ni traslapos y solo en el perímetro, relleno en m$^2$ sobre la huella, bajadas que repiten la cifra de canales y cerámica en todo el interior cuando las EETT la piden en cocina y baño.
\item \textbf{Efecto}: la escalerilla queda sobrepresupuestada en unos \$\vSob{} mil y la cadena (hormigón y acero) subpresupuestada en unos \$\vFal{} mil. El capítulo 3 parece razonable en total, pero la plata está mal repartida y falta justo en una partida estructural.
\end{itemize}
\end{resultado}

\section{Fuentes y método}
Fuentes públicas del FSPR (2023): Anexo 5.3, Planimetría Sede Social Tipo C (lámina 1/1); Anexo 3.3, Presupuesto; Anexo 4.3, EETT. Criterios de medición de NCh353.Of2000 (texto de la norma); ordenación de partidas según NCh1156; masas nominales y largos de barra según NCh204:2020; requisitos de albañilería armada según NCh1928.Of1993 Mod. 2009 (cláusulas 6, 7 y Anexo A); largos de desarrollo y empalmes según ACI 318-19, capítulo 25.

\begin{juntos}
La planta 1:50 se extrajo del PDF vectorial a DXF con un script propio y se escaló con la cota total de \SI{12.85}{m}. La escala se comprobó con cotas que no se usaron para calibrar:
\begin{table}
\centering\small
\begin{tabular}{lSSr}
\toprule
Cota del plano & {Cota [m]} & {Medido [m]} & Diferencia\\
\midrule
Proyección de techumbre & 6.50 & 6.502 & 0.0\,\%\\
Muro exterior a muro exterior & 5.30 & 5.301 & 0.0\,\%\\
Espesor de muro (3 muros) & 0.15 & {0.148--0.151} & $-1.2$ a $+0.6$\,\%\\
\bottomrule
\end{tabular}
\caption{Verificación de la escala de la geometría extraída.}
\label{tab:escala}
\end{table}
\end{juntos}

El largo de muros, base de casi todas las partidas, se obtuvo por tres caminos independientes: ejes del DXF (\SI{\vLeje}{m}), recubicación ciega hecha por otro analista desde el PDF sin conocer este resultado (\SI{37.32}{m}) y área rasterizada de muros dividida por el espesor (\SI{37.31}{m}). La dispersión es de 0.4\,\%.

\begin{juntos}
\begin{figure}
\centering
\begin{tikzpicture}[scale=1.05]
\draw[muro] (0.000,-1.951)--(0.000,-0.148);
\draw[muro] (0.000,-0.148)--(-1.151,-0.148);
\draw[muro] (-1.151,-0.148)--(-1.151,0.000);
\draw[muro] (-1.151,0.000)--(0.000,0.000);
\draw[muro] (0.151,0.000)--(1.002,0.000);
\draw[muro] (1.401,0.000)--(3.452,0.000);
\draw[muro] (6.303,-0.148)--(7.653,-0.148);
\draw[muro] (6.303,0.000)--(6.904,0.000);
\draw[muro] (6.904,0.000)--(6.904,0.501);
\draw[muro] (6.904,0.501)--(7.052,0.501);
\draw[muro] (7.052,0.501)--(7.052,0.000);
\draw[muro] (10.356,0.000)--(10.356,0.501);
\draw[muro] (10.356,0.501)--(10.504,0.501);
\draw[muro] (10.504,0.501)--(10.504,0.000);
\draw[muro] (10.504,0.000)--(10.504,-5.301);
\draw[muro] (0.151,-1.951)--(0.151,-0.148);
\draw[muro] (0.151,-0.148)--(1.002,-0.148);
\draw[muro] (1.401,-0.148)--(4.201,-0.148);
\draw[muro] (3.603,0.000)--(4.201,0.000);
\draw[muro] (0.000,-3.250)--(0.151,-3.250);
\draw[muro] (4.802,-0.148)--(4.802,0.000);
\draw[muro] (4.201,-0.148)--(4.201,0.000);
\draw[muro] (0.000,-1.951)--(0.151,-1.951);
\draw[muro] (3.076,-5.150)--(4.201,-5.150);
\draw[muro] (1.401,-5.301)--(3.452,-5.301);
\draw[muro] (-0.002,-3.250)--(-0.002,-5.150);
\draw[muro] (-0.002,-5.150)--(-1.151,-5.150);
\draw[muro] (-1.151,-5.150)--(-1.151,-5.301);
\draw[muro] (-1.151,-5.301)--(0.000,-5.301);
\draw[muro] (0.151,-5.301)--(1.002,-5.301);
\draw[muro] (0.151,-3.250)--(0.151,-5.150);
\draw[muro] (0.151,-5.150)--(1.002,-5.150);
\draw[muro] (1.401,-5.150)--(3.352,-5.150);
\draw[muro] (10.356,-5.152)--(10.356,-0.148);
\draw[muro] (1.401,-5.150)--(1.401,-5.301);
\draw[muro] (1.002,-5.150)--(1.002,-5.301);
\draw[muro] (0.151,-3.449)--(2.352,-3.449);
\draw[muro] (1.401,-0.150)--(1.401,0.000);
\draw[muro] (1.002,-0.148)--(1.002,0.000);
\draw[muro] (3.603,0.000)--(3.603,0.501);
\draw[muro] (3.603,0.501)--(3.452,0.501);
\draw[muro] (3.452,0.501)--(3.452,0.000);
\draw[muro] (0.151,0.000)--(0.151,0.501);
\draw[muro] (0.151,0.501)--(0.000,0.501);
\draw[muro] (0.000,0.501)--(0.000,0.000);
\draw[muro] (3.603,-5.301)--(3.603,-5.802);
\draw[muro] (3.603,-5.802)--(3.452,-5.802);
\draw[muro] (3.452,-5.802)--(3.452,-5.301);
\draw[muro] (0.151,-5.301)--(0.151,-5.802);
\draw[muro] (0.151,-5.802)--(0.000,-5.802);
\draw[muro] (0.000,-5.802)--(0.000,-5.301);
\draw[muro] (4.201,-5.301)--(3.603,-5.301);
\draw[muro] (6.303,-0.148)--(6.303,0.000);
\draw[muro] (5.702,-0.148)--(5.702,0.000);
\draw[muro] (5.702,-0.148)--(4.802,-0.148);
\draw[muro] (5.702,0.000)--(4.802,0.000);
\draw[muro] (9.154,-0.148)--(9.154,0.000);
\draw[muro] (9.755,-0.148)--(9.755,0.000);
\draw[muro] (7.653,-0.148)--(7.653,0.000);
\draw[muro] (8.254,-0.148)--(8.254,0.000);
\draw[muro] (10.356,-0.148)--(9.755,-0.148);
\draw[muro] (10.356,0.000)--(9.755,0.000);
\draw[muro] (9.154,0.000)--(8.254,0.000);
\draw[muro] (9.154,-0.148)--(8.254,-0.148);
\draw[muro] (7.653,0.000)--(7.052,0.000);
\draw[muro] (6.303,-5.152)--(7.653,-5.152);
\draw[muro] (6.303,-5.301)--(6.904,-5.301);
\draw[muro] (6.904,-5.301)--(6.904,-5.802);
\draw[muro] (6.904,-5.802)--(7.052,-5.802);
\draw[muro] (7.052,-5.802)--(7.052,-5.301);
\draw[muro] (10.356,-5.301)--(10.356,-5.802);
\draw[muro] (10.356,-5.802)--(10.504,-5.802);
\draw[muro] (10.504,-5.802)--(10.504,-5.301);
\draw[muro] (4.802,-5.152)--(4.802,-5.301);
\draw[muro] (4.201,-5.152)--(4.201,-5.301);
\draw[muro] (6.303,-5.152)--(6.303,-5.301);
\draw[muro] (5.702,-5.152)--(5.702,-5.301);
\draw[muro] (5.702,-5.152)--(4.802,-5.152);
\draw[muro] (5.702,-5.301)--(4.802,-5.301);
\draw[muro] (9.154,-5.152)--(9.154,-5.301);
\draw[muro] (9.755,-5.152)--(9.755,-5.301);
\draw[muro] (7.653,-5.152)--(7.653,-5.301);
\draw[muro] (8.254,-5.152)--(8.254,-5.301);
\draw[muro] (10.356,-5.152)--(9.755,-5.152);
\draw[muro] (10.356,-5.301)--(9.755,-5.301);
\draw[muro] (9.154,-5.301)--(8.254,-5.301);
\draw[muro] (9.154,-5.152)--(8.254,-5.152);
\draw[muro] (7.653,-5.301)--(7.052,-5.301);
\draw[elem] (0.151,-3.250)--(0.151,-3.201);
\draw[elem] (0.151,-3.201)--(0.102,-3.201);
\draw[elem] (0.102,-3.201)--(0.102,-3.250);
\draw[elem] (-0.575,-3.150)--(-0.575,-3.201);
\draw[elem] (-0.425,-3.150)--(-0.425,-3.201);
\draw[elem] (4.201,-0.148)--(4.802,-0.148);
\draw[elem] (4.753,-0.099)--(4.252,-0.099);
\draw[elem] (4.201,-0.099)--(4.252,-0.099);
\draw[elem] (4.252,-0.099)--(4.252,0.000);
\draw[elem] (4.252,0.000)--(4.201,0.000);
\draw[elem] (4.802,-0.099)--(4.753,-0.099);
\draw[elem] (4.753,-0.099)--(4.753,0.000);
\draw[elem] (4.753,0.000)--(4.802,0.000);
\draw[elem] (4.753,0.000)--(4.252,0.000);
\draw[elem] (0.151,-1.951)--(0.151,-1.999);
\draw[elem] (0.151,-1.999)--(0.102,-1.999);
\draw[elem] (0.102,-1.999)--(0.102,-1.951);
\draw[elem] (2.397,-5.150)--(2.397,-4.551);
\draw[elem] (0.374,-4.575)--(0.654,-4.575);
\draw[elem] (2.397,-4.451)--(2.297,-4.451);
\draw[elem] (2.297,-4.451)--(2.297,-5.150);
\draw[elem] (0.374,-4.935)--(0.654,-4.935);
\draw[elem] (0.174,-4.861)--(0.174,-4.651);
\draw[elem] (0.355,-4.861)--(0.355,-4.651);
\draw[elem] (0.556,-4.615)--(0.654,-4.615);
\draw[elem] (0.556,-4.895)--(0.654,-4.895);
\draw[elem] (2.214,-5.131)--(2.214,-4.965);
\draw[elem] (1.543,-4.965)--(1.543,-5.131);
\draw[elem] (1.563,-5.150)--(2.193,-5.150);
\draw[elem] (1.582,-4.969)--(1.582,-4.965);
\draw[elem] (2.176,-4.969)--(2.176,-4.965);
\draw[elem] (2.117,-5.031)--(1.641,-5.031);
\draw[elem] (1.352,-5.201)--(1.051,-5.201);
\draw[elem] (1.401,-5.201)--(1.352,-5.201);
\draw[elem] (1.352,-5.201)--(1.352,-5.301);
\draw[elem] (1.352,-5.301)--(1.401,-5.301);
\draw[elem] (1.352,-5.301)--(1.051,-5.301);
\draw[elem] (1.002,-5.201)--(1.051,-5.201);
\draw[elem] (1.051,-5.201)--(1.051,-5.301);
\draw[elem] (1.051,-5.301)--(1.002,-5.301);
\draw[elem] (0.853,-5.071)--(0.293,-5.071);
\draw[elem] (0.853,-5.093)--(0.293,-5.093);
\draw[elem] (0.928,-4.366)--(0.229,-4.366);
\draw[elem] (0.928,-4.343)--(0.229,-4.343);
\draw[elem] (0.151,-4.328)--(0.229,-4.337);
\draw[elem] (0.229,-4.337)--(0.229,-4.356);
\draw[elem] (0.229,-4.356)--(0.229,-4.375);
\draw[elem] (0.229,-4.375)--(0.151,-4.383);
\draw[elem] (0.151,-3.449)--(2.427,-3.449);
\draw[elem] (2.427,-3.449)--(2.427,-3.350);
\draw[elem] (2.427,-3.350)--(0.151,-3.350);
\draw[elem] (0.187,-1.632)--(0.187,-0.851);
\draw[elem] (0.187,-0.851)--(0.189,-0.845);
\draw[elem] (0.189,-0.845)--(0.195,-0.840);
\draw[elem] (0.195,-1.641)--(0.189,-1.639);
\draw[elem] (0.189,-1.639)--(0.187,-1.630);
\draw[elem] (0.197,-1.641)--(0.611,-1.641);
\draw[elem] (0.611,-1.641)--(0.618,-1.639);
\draw[elem] (0.618,-1.639)--(0.620,-1.630);
\draw[elem] (0.622,-1.632)--(0.622,-0.851);
\draw[elem] (0.195,-0.840)--(0.611,-0.840);
\draw[elem] (0.611,-0.840)--(0.618,-0.845);
\draw[elem] (0.618,-0.845)--(0.620,-0.851);
\draw[elem] (0.242,-1.619)--(0.227,-1.617);
\draw[elem] (0.227,-1.617)--(0.217,-1.611);
\draw[elem] (0.217,-1.611)--(0.208,-1.600);
\draw[elem] (0.208,-1.600)--(0.206,-1.585);
\draw[elem] (0.206,-1.585)--(0.206,-0.896);
\draw[elem] (0.206,-0.896)--(0.208,-0.883);
\draw[elem] (0.208,-0.883)--(0.217,-0.870);
\draw[elem] (0.217,-0.870)--(0.227,-0.864);
\draw[elem] (0.227,-0.864)--(0.240,-0.862);
\draw[elem] (0.242,-1.622)--(0.567,-1.622);
\draw[elem] (0.242,-0.862)--(0.567,-0.862);
\draw[elem] (0.567,-0.862)--(0.580,-0.864);
\draw[elem] (0.580,-0.864)--(0.590,-0.870);
\draw[elem] (0.590,-0.870)--(0.599,-0.883);
\draw[elem] (0.599,-0.883)--(0.601,-0.896);
\draw[elem] (0.601,-0.896)--(0.601,-1.585);
\draw[elem] (0.601,-1.585)--(0.599,-1.600);
\draw[elem] (0.599,-1.600)--(0.590,-1.611);
\draw[elem] (0.590,-1.611)--(0.580,-1.617);
\draw[elem] (0.580,-1.617)--(0.567,-1.619);
\draw[elem] (0.348,-1.023)--(0.357,-1.023);
\draw[elem] (0.359,-1.023)--(0.365,-1.023);
\draw[elem] (0.369,-1.023)--(0.376,-1.023);
\draw[elem] (0.386,-1.023)--(0.380,-1.023);
\draw[elem] (0.393,-1.023)--(0.399,-1.023);
\draw[elem] (0.403,-1.023)--(0.410,-1.023);
\draw[elem] (0.420,-1.023)--(0.412,-1.023);
\draw[elem] (0.431,-1.023)--(0.422,-1.023);
\draw[elem] (0.435,-1.023)--(0.442,-1.023);
\draw[elem] (0.446,-1.023)--(0.452,-1.023);
\draw[elem] (0.456,-1.023)--(0.463,-1.023);
\draw[elem] (0.380,-1.021)--(0.382,-1.012);
\draw[elem] (0.382,-1.012)--(0.386,-1.006);
\draw[elem] (0.386,-1.006)--(0.395,-1.002);
\draw[elem] (0.395,-1.002)--(0.403,-1.000);
\draw[elem] (0.403,-1.000)--(0.412,-1.002);
\draw[elem] (0.412,-1.002)--(0.418,-1.006);
\draw[elem] (0.418,-1.006)--(0.422,-1.012);
\draw[elem] (0.422,-1.012)--(0.425,-1.021);
\draw[elem] (0.425,-1.021)--(0.422,-1.031);
\draw[elem] (0.422,-1.031)--(0.418,-1.038);
\draw[elem] (0.418,-1.038)--(0.412,-1.042);
\draw[elem] (0.412,-1.042)--(0.403,-1.044);
\draw[elem] (0.403,-1.044)--(0.401,-1.044);
\draw[elem] (0.401,-1.044)--(0.391,-1.040);
\draw[elem] (0.391,-1.040)--(0.384,-1.031);
\draw[elem] (0.384,-1.031)--(0.380,-1.021);
\draw[elem] (0.261,-1.326)--(0.263,-1.598);
\draw[elem] (0.484,-0.893)--(0.501,-0.898);
\draw[elem] (0.501,-0.898)--(0.516,-0.904);
\draw[elem] (0.516,-0.904)--(0.531,-0.913);
\draw[elem] (0.531,-0.913)--(0.544,-0.923);
\draw[elem] (0.544,-0.923)--(0.552,-0.936);
\draw[elem] (0.552,-0.936)--(0.560,-0.949);
\draw[elem] (0.560,-0.949)--(0.567,-0.966);
\draw[elem] (0.567,-0.966)--(0.569,-0.983);
\draw[elem] (0.571,-0.983)--(0.573,-1.034);
\draw[elem] (0.573,-1.034)--(0.575,-1.084);
\draw[elem] (0.575,-1.084)--(0.573,-1.135);
\draw[elem] (0.573,-1.135)--(0.567,-1.186);
\draw[elem] (0.567,-1.186)--(0.565,-1.199);
\draw[elem] (0.565,-1.199)--(0.558,-1.212);
\draw[elem] (0.558,-1.212)--(0.552,-1.225);
\draw[elem] (0.552,-1.225)--(0.544,-1.235);
\draw[elem] (0.544,-1.235)--(0.533,-1.244);
\draw[elem] (0.533,-1.244)--(0.522,-1.250);
\draw[elem] (0.522,-1.250)--(0.510,-1.256);
\draw[elem] (0.510,-1.256)--(0.495,-1.259);
\draw[elem] (0.495,-1.259)--(0.450,-1.265);
\draw[elem] (0.450,-1.265)--(0.403,-1.265);
\draw[elem] (0.403,-1.265)--(0.357,-1.265);
\draw[elem] (0.357,-1.265)--(0.312,-1.259);
\draw[elem] (0.312,-1.261)--(0.297,-1.256);
\draw[elem] (0.297,-1.256)--(0.285,-1.252);
\draw[elem] (0.285,-1.252)--(0.274,-1.244);
\draw[elem] (0.274,-1.244)--(0.263,-1.235);
\draw[elem] (0.263,-1.235)--(0.255,-1.225);
\draw[elem] (0.255,-1.225)--(0.248,-1.212);
\draw[elem] (0.248,-1.212)--(0.242,-1.199);
\draw[elem] (0.242,-1.199)--(0.240,-1.186);
\draw[elem] (0.240,-1.186)--(0.234,-1.135);
\draw[elem] (0.234,-1.135)--(0.231,-1.084);
\draw[elem] (0.231,-1.084)--(0.234,-1.034);
\draw[elem] (0.234,-1.034)--(0.236,-0.983);
\draw[elem] (0.236,-0.983)--(0.240,-0.966);
\draw[elem] (0.240,-0.966)--(0.244,-0.949);
\draw[elem] (0.244,-0.949)--(0.253,-0.936);
\draw[elem] (0.253,-0.936)--(0.263,-0.923);
\draw[elem] (0.263,-0.923)--(0.276,-0.913);
\draw[elem] (0.276,-0.913)--(0.289,-0.904);
\draw[elem] (0.289,-0.904)--(0.306,-0.898);
\draw[elem] (0.306,-0.898)--(0.321,-0.893);
\draw[elem] (0.321,-0.893)--(0.363,-0.889);
\draw[elem] (0.363,-0.889)--(0.403,-0.887);
\draw[elem] (0.403,-0.887)--(0.444,-0.889);
\draw[elem] (0.444,-0.889)--(0.484,-0.893);
\draw[elem] (0.333,-1.237)--(0.321,-1.235);
\draw[elem] (0.321,-1.235)--(0.308,-1.229);
\draw[elem] (0.308,-1.229)--(0.289,-1.214);
\draw[elem] (0.289,-1.214)--(0.274,-1.193);
\draw[elem] (0.274,-1.193)--(0.270,-1.182);
\draw[elem] (0.270,-1.182)--(0.265,-1.167);
\draw[elem] (0.265,-1.167)--(0.263,-1.123);
\draw[elem] (0.263,-1.123)--(0.261,-1.078);
\draw[elem] (0.261,-1.078)--(0.263,-1.034);
\draw[elem] (0.263,-1.034)--(0.268,-0.991);
\draw[elem] (0.268,-0.991)--(0.274,-0.972);
\draw[elem] (0.274,-0.972)--(0.282,-0.957);
\draw[elem] (0.282,-0.957)--(0.297,-0.946);
\draw[elem] (0.297,-0.946)--(0.314,-0.940);
\draw[elem] (0.473,-1.239)--(0.437,-1.242);
\draw[elem] (0.437,-1.242)--(0.403,-1.244);
\draw[elem] (0.403,-1.244)--(0.369,-1.242);
\draw[elem] (0.369,-1.242)--(0.333,-1.239);
\draw[elem] (0.314,-0.940)--(0.359,-0.934);
\draw[elem] (0.359,-0.934)--(0.403,-0.932);
\draw[elem] (0.403,-0.932)--(0.448,-0.934);
\draw[elem] (0.448,-0.934)--(0.493,-0.940);
\draw[elem] (0.493,-0.940)--(0.510,-0.946);
\draw[elem] (0.510,-0.946)--(0.524,-0.957);
\draw[elem] (0.524,-0.957)--(0.533,-0.972);
\draw[elem] (0.533,-0.972)--(0.539,-0.991);
\draw[elem] (0.539,-0.991)--(0.544,-1.034);
\draw[elem] (0.544,-1.034)--(0.546,-1.078);
\draw[elem] (0.546,-1.078)--(0.546,-1.123);
\draw[elem] (0.546,-1.123)--(0.541,-1.167);
\draw[elem] (0.541,-1.167)--(0.537,-1.180);
\draw[elem] (0.537,-1.180)--(0.533,-1.193);
\draw[elem] (0.533,-1.193)--(0.518,-1.214);
\draw[elem] (0.518,-1.214)--(0.499,-1.229);
\draw[elem] (0.499,-1.229)--(0.486,-1.235);
\draw[elem] (0.486,-1.235)--(0.473,-1.237);
\draw[elem] (0.295,-1.301)--(0.295,-1.598);
\draw[elem] (0.329,-1.284)--(0.329,-1.598);
\draw[elem] (0.361,-1.284)--(0.361,-1.598);
\draw[elem] (0.395,-1.284)--(0.395,-1.598);
\draw[elem] (0.529,-1.318)--(0.529,-1.598);
\draw[elem] (0.495,-1.284)--(0.495,-1.598);
\draw[elem] (0.463,-1.284)--(0.463,-1.598);
\draw[elem] (0.429,-1.284)--(0.429,-1.598);
\draw[elem] (0.403,-1.070)--(0.403,-1.061);
\draw[elem] (0.403,-1.051)--(0.403,-1.059);
\draw[elem] (0.403,-1.046)--(0.403,-1.040);
\draw[elem] (0.403,-1.036)--(0.403,-1.029);
\draw[elem] (0.403,-1.025)--(0.403,-1.019);
\draw[elem] (0.403,-1.014)--(0.403,-1.008);
\draw[elem] (0.403,-1.004)--(0.403,-0.995);
\draw[elem] (0.403,-0.993)--(0.403,-0.987);
\draw[elem] (0.403,-0.983)--(0.403,-0.974);
\draw[elem] (0.652,-1.749)--(0.652,-0.649);
\draw[elem] (3.947,-5.150)--(3.947,-4.551);
\draw[elem] (3.947,-4.451)--(3.947,-3.350);
\draw[elem] (3.947,-3.350)--(4.048,-3.350);
\draw[elem] (4.048,-3.350)--(4.048,-5.150);
\draw[elem] (2.323,-4.451)--(2.323,-4.400);
\draw[elem] (2.323,-4.400)--(2.371,-4.400);
\draw[elem] (2.371,-4.400)--(2.371,-4.451);
\draw[elem] (2.577,-4.481)--(3.172,-4.481);
\draw[elem] (2.647,-4.500)--(3.698,-4.500);
\draw[elem] (2.577,-4.500)--(3.777,-4.500);
\draw[elem] (3.172,-4.481)--(3.172,-4.521);
\draw[elem] (3.172,-4.521)--(3.777,-4.521);
\draw[elem] (3.947,-4.451)--(3.777,-4.451);
\draw[elem] (3.777,-4.451)--(3.777,-4.551);
\draw[elem] (3.777,-4.551)--(3.947,-4.551);
\draw[elem] (2.371,-3.449)--(2.371,-3.500);
\draw[elem] (2.371,-3.500)--(2.323,-3.500);
\draw[elem] (2.323,-3.500)--(2.323,-3.449);
\draw[elem] (2.397,-4.451)--(2.577,-4.451);
\draw[elem] (2.577,-4.451)--(2.577,-4.551);
\draw[elem] (2.577,-4.551)--(2.397,-4.551);
\draw[elem] (2.677,-0.148)--(2.677,-0.649);
\draw[elem] (2.677,-1.649)--(2.677,-1.851);
\draw[elem] (2.677,-1.851)--(2.577,-1.851);
\draw[elem] (2.577,-1.851)--(2.577,-1.649);
\draw[elem] (2.579,-0.649)--(2.581,-0.148);
\draw[elem] (0.151,-1.851)--(0.955,-1.851);
\draw[elem] (0.955,-1.851)--(0.955,-1.749);
\draw[elem] (0.955,-1.749)--(0.151,-1.749);
\draw[elem] (0.652,-0.649)--(1.652,-0.649);
\draw[elem] (1.652,-0.649)--(1.652,-0.148);
\draw[elem] (-1.278,-0.099)--(-1.278,-0.048);
\draw[elem] (-1.278,-1.600)--(-1.278,-1.549);
\draw[elem] (1.352,-0.099)--(1.051,-0.099);
\draw[elem] (1.401,-0.099)--(1.352,-0.099);
\draw[elem] (1.352,-0.099)--(1.352,0.000);
\draw[elem] (1.352,0.000)--(1.401,0.000);
\draw[elem] (1.401,-5.150)--(1.002,-5.150);
\draw[elem] (1.352,0.000)--(1.051,0.000);
\draw[elem] (1.401,-0.148)--(1.002,-0.148);
\draw[elem] (1.002,-0.099)--(1.051,-0.099);
\draw[elem] (1.051,-0.099)--(1.051,0.000);
\draw[elem] (1.051,0.000)--(1.002,0.000);
\draw[elem] (2.677,-1.649)--(2.577,-1.649);
\draw[elem] (2.677,-0.649)--(2.579,-0.649);
\draw[elem] (5.702,-0.148)--(6.303,-0.148);
\draw[elem] (6.252,-0.099)--(5.753,-0.099);
\draw[elem] (5.702,-0.099)--(5.753,-0.099);
\draw[elem] (5.753,-0.099)--(5.753,0.000);
\draw[elem] (5.753,0.000)--(5.702,0.000);
\draw[elem] (6.303,-0.099)--(6.252,-0.099);
\draw[elem] (6.252,-0.099)--(6.252,0.000);
\draw[elem] (6.252,0.000)--(6.303,0.000);
\draw[elem] (6.252,0.000)--(5.753,0.000);
\draw[elem] (9.755,-0.148)--(9.154,-0.148);
\draw[elem] (9.205,-0.099)--(9.704,-0.099);
\draw[elem] (9.755,-0.099)--(9.704,-0.099);
\draw[elem] (9.704,-0.099)--(9.704,0.000);
\draw[elem] (9.704,0.000)--(9.755,0.000);
\draw[elem] (9.154,-0.099)--(9.205,-0.099);
\draw[elem] (9.205,-0.099)--(9.205,0.000);
\draw[elem] (9.205,0.000)--(9.154,0.000);
\draw[elem] (9.205,0.000)--(9.704,0.000);
\draw[elem] (8.254,-0.148)--(7.653,-0.148);
\draw[elem] (7.704,-0.099)--(8.203,-0.099);
\draw[elem] (8.254,-0.099)--(8.203,-0.099);
\draw[elem] (8.203,-0.099)--(8.203,0.000);
\draw[elem] (8.203,0.000)--(8.254,0.000);
\draw[elem] (7.653,-0.099)--(7.704,-0.099);
\draw[elem] (7.704,-0.099)--(7.704,0.000);
\draw[elem] (7.704,0.000)--(7.653,0.000);
\draw[elem] (7.704,0.000)--(8.203,0.000);
\draw[elem] (4.201,-5.152)--(4.802,-5.152);
\draw[elem] (4.753,-5.203)--(4.252,-5.203);
\draw[elem] (4.201,-5.203)--(4.252,-5.203);
\draw[elem] (4.252,-5.203)--(4.252,-5.301);
\draw[elem] (4.252,-5.301)--(4.201,-5.301);
\draw[elem] (4.802,-5.203)--(4.753,-5.203);
\draw[elem] (4.753,-5.203)--(4.753,-5.301);
\draw[elem] (4.753,-5.301)--(4.802,-5.301);
\draw[elem] (4.753,-5.301)--(4.252,-5.301);
\draw[elem] (5.702,-5.152)--(6.303,-5.152);
\draw[elem] (6.252,-5.201)--(5.753,-5.201);
\draw[elem] (5.702,-5.201)--(5.753,-5.201);
\draw[elem] (5.753,-5.201)--(5.753,-5.301);
\draw[elem] (5.753,-5.301)--(5.702,-5.301);
\draw[elem] (6.303,-5.201)--(6.252,-5.201);
\draw[elem] (6.252,-5.201)--(6.252,-5.301);
\draw[elem] (6.252,-5.301)--(6.303,-5.301);
\draw[elem] (6.252,-5.301)--(5.753,-5.301);
\draw[elem] (9.755,-5.152)--(9.154,-5.152);
\draw[elem] (9.205,-5.201)--(9.704,-5.201);
\draw[elem] (9.755,-5.201)--(9.704,-5.201);
\draw[elem] (9.704,-5.201)--(9.704,-5.301);
\draw[elem] (9.704,-5.301)--(9.755,-5.301);
\draw[elem] (9.154,-5.201)--(9.205,-5.201);
\draw[elem] (9.205,-5.201)--(9.205,-5.301);
\draw[elem] (9.205,-5.301)--(9.154,-5.301);
\draw[elem] (9.205,-5.301)--(9.704,-5.301);
\draw[elem] (8.254,-5.152)--(7.653,-5.152);
\draw[elem] (7.704,-5.201)--(8.203,-5.201);
\draw[elem] (8.254,-5.201)--(8.203,-5.201);
\draw[elem] (8.203,-5.201)--(8.203,-5.301);
\draw[elem] (8.203,-5.301)--(8.254,-5.301);
\draw[elem] (7.653,-5.201)--(7.704,-5.201);
\draw[elem] (7.704,-5.201)--(7.704,-5.301);
\draw[elem] (7.704,-5.301)--(7.653,-5.301);
\draw[elem] (7.704,-5.301)--(8.203,-5.301);
\coordinate (V0) at (-1.151,0); \coordinate (P0) at (0,0); \coordinate (P1) at (10.504,0);
\coordinate (Q1) at (10.504,-5.301);
\cota[7mm]{V0}{P0}{\num{1.15}}
\cota[7mm]{P0}{P1}{\num{10.50}}
\cota[-9mm]{P1}{Q1}{\num{5.30}}
\node[etq,align=center] at (5.2,-2.65) {Salón multiuso};
\node[etq,align=left,anchor=west] at (-1.10,-2.65) {Vestíbulo};
\end{tikzpicture}
\caption{Planta de muros extraída del plano (muros cortados en negro, ventanas y artefactos en azul). Cotas medidas en el DXF, en metros; acceso a la izquierda.}
\label{fig:planta}
\end{figure}
\end{juntos}

La verificación tuvo tres capas: la planilla se contrasta fila por fila con un script de Python que calcula por otro camino (22 de 22 controles en OK); un revisor externo auditó el código y los datos contra las EETT; y un segundo analista recubicó a ciegas las partidas con mayor diferencia. Base de cada cantidad: M, medida en el plano; E, con sección o armadura de las EETT; I, con sección inferida porque las EETT remiten a planos de estructuras que no están en el expediente; S, con un supuesto declarado en la sección~\ref{sec:supuestos}.

\section{Comparación partida por partida}
\begin{juntos}
\begin{table}
\centering\small
\begin{tabular}{llcSSrcl}
\toprule
Ítem & Partida & Unidad & {Presupuesto} & {Cubicación} & {Dif. [\%]} & Base & Causa\\
\midrule
2.1 & Nivelación y trazado & m$^2$ & \num{62.70} & \num{62.38} & +0.5 & M & \\
2.2 & Escarpe e = 0.15 m & m$^3$ & \num{9.40} & \num{9.36} & +0.5 & M & \\
2.3 & Excavación de fundaciones & m$^3$ & \num{10.25} & \num{9.74} & +5.2 & I & Criterio\\
2.4 & Relleno estabilizado e = 0.15 m & m$^3$ & \num{9.40} & \num{7.65} & +22.8 & M & Error ppto.\\
2.5 & Emplantillado H-5 & m$^3$ & \num{0.76} & \num{0.75} & +1.4 & I & \\
2.6 & Fundaciones H-25 & m$^3$ & \num{9.20} & \num{8.99} & +2.3 & I & \\
2.7 & Sobrecimientos & m$^3$ & \num{1.15} & \num{1.12} & +2.3 & I & \\
2.8 & Radier e = 0.08 m & m$^3$ & \num{4.26} & \num{4.08} & +4.3 & M & \\
2.9 & Moldaje sobrecimientos & m$^2$ & \num{16.00} & \num{14.99} & +6.8 & I & Criterio\\
3.1 & Ladrillo gran titán & m$^2$ & \num{90.00} & \num{88.56} & +1.6 & M & Criterio\\
3.2 & Escalerilla & kg & \num{186.50} & \num{67.12} & +177.8 & S & Error ppto.\\
3.4 & Hormigón cadena & m$^3$ & \num{1.28} & \num{1.69} & -24.1 & E & Error ppto.\\
3.6 & Moldaje cadena & m$^2$ & \num{30.00} & \num{23.77} & +26.2 & E & Criterio\\
3.5 & Enfierradura cadena & kg & \num{93.00} & \num{194.77} & -52.3 & E & Error ppto.\\
3.3 & Tensores & kg & \num{174.00} & \num{136.27} & +27.7 & M/S & Criterio\\
5.5 & Cubierta OSB 11 mm & m$^2$ & \num{87.50} & \num{89.90} & -2.7 & M & \\
7.1 & Canales & ml & \num{26.00} & \num{25.70} & +1.2 & M & \\
8.1 & Bajadas AALL & ml & \num{26.00} & \num{15.60} & +66.7 & S & Error ppto.\\
10.3 & Cielo yeso-cartón & m$^2$ & \num{53.30} & \num{51.03} & +4.4 & M & \\
11.1 & Cerámica 30 x 30 & m$^2$ & \num{51.00} & \num{7.95} & +541.5 & M & Error ppto.\\
\bottomrule
\end{tabular}
\caption{Presupuesto tipo frente a la cubicación independiente (v6). Diferencia positiva: el presupuesto cubica más.}
\label{tab:comp}
\end{table}

Las partidas sin causa indicada están dentro de $\pm5$\,\%. El ladrillo se incluye en la conciliación porque su diferencia cambió de $+15$\,\% a $+1.6$\,\% al aplicar NCh 353.
\end{juntos}

\section{Por qué difiere el presupuesto}
\begin{juntos}
Para cada partida fuera de tolerancia se buscó el criterio que reproduce la cifra del presupuesto. Se clasifica como criterio cuando la cifra sale de una práctica legítima (NCh 353 o práctica de obra) y como error cuando solo se explica con una cantidad que no corresponde a la geometría ni a las EETT, una unidad equivocada o una cifra repetida.
\begin{table}
\centering\small
\begin{tabularx}{\linewidth}{lp{2.9cm}lXS[table-format=3.2]S[table-format=3.2]r}
\toprule
Ítem & Partida & Causa & Criterio que reproduce la cifra & {Reconstr.} & {Ppto.} & {Dif. [\%]}\\
\midrule
3.1 & Ladrillo gran titán & Criterio & Sin descontar vanos: $L\times2.40$ & \num{89.93} & \num{90.00} & -0.1\\
2.3 & Excavación de fundaciones & Criterio & Sobreexcavación de 5\,\% & \num{10.23} & \num{10.25} & -0.2\\
2.9 & Moldaje sobrecimientos & Criterio & Moldaje neto $\times1.07$ & \num{16.04} & \num{16.00} & +0.2\\
3.3 & Tensores & Criterio & Todos $\phi$10: 85 barras de \SI{3.30}{m} & \num{173.07} & \num{174.00} & -0.5\\
3.6 & Moldaje cadena & Criterio & Tablero de \SI{0.40}{m} por cara & \num{29.98} & \num{30.00} & -0.1\\
5.5 & Cubierta OSB 11 mm & Criterio & Vertientes de \SI{12.50}{m}$\times$\SI{3.50}{m} & \num{87.50} & \num{87.50} & +0.0\\
2.4 & Relleno estabilizado e = 0.15 m & Error ppto. & Huella completa $\times0.15$ en m$^2$ & \num{9.36} & \num{9.40} & -0.5\\
3.2 & Escalerilla & Error ppto. & 10 filas $\times$ \SI{0.50}{kg/m}: doble de la masa real & \num{187.35} & \num{186.50} & +0.5\\
3.4 & Hormigón cadena & Error ppto. & Cadena solo en el perímetro, muro de \SI{0.14}{m} & \num{1.30} & \num{1.28} & +1.8\\
3.5 & Enfierradura cadena & Error ppto. & 4\,$\phi$10 sin estribos ni traslapos & \num{92.48} & \num{93.00} & -0.6\\
8.1 & Bajadas AALL & Error ppto. & Repite la cifra de canales & \num{26.00} & \num{26.00} & +0.0\\
11.1 & Cerámica 30 x 30 & Error ppto. & Cerámica en todo el interior & \num{51.03} & \num{51.00} & +0.1\\
\bottomrule
\end{tabularx}
\caption{Reconstrucción de las cifras del presupuesto: todas salvo el cielo (\SI{53.3}{m^2}, sin resolver) se reproducen dentro de $\pm2$\,\%.}
\label{tab:conc}
\end{table}

\end{juntos}

\begin{itemize}[nosep]
\item \textbf{Ladrillo.} NCh 353 no descuenta vanos menores a \SI{1.5}{m^2} en ladrillo hecho a máquina y descuenta al 50\,\% los de 1.5 a \SI{3}{m^2}: las ventanas (\SI{1.08}{m^2}) no se descuentan y la puerta (\SI{2.73}{m^2}) a medias. Resultado: \SI{\vLad}{m^2}, a 1.6\,\% del presupuesto. Si los vanos llevaran pilar, \SI{\vLadPilar}{m^2}.
\item \textbf{Cadena.} Con 6\,$\phi$10 y estribos $\phi$6 a 0.20 (EETT 3.4), traslapos de 40\,$d$ (EETT 3.5) y +5\,\% de patas (NCh 353): \SI{\vFcad}{kg}, el doble de los \SI{93}{kg} presupuestados. El hormigón presupuestado (\SI{1.28}{m^3}) corresponde a una cadena solo en el perímetro, pero el moldaje (\SI{30}{m^2}) corresponde al largo completo: el presupuesto es incoherente consigo mismo.
\item \textbf{Escalerilla.} Una fila cada 3 hiladas (EETT 3.2) da 9 filas bajo la cadena y \SI{\vLesc}{m} netos de vanos; con alambre de \SI{4}{mm} (\SI{\vEscKgm}{kg/m}) son \SI{\vEsc}{kg}. El presupuesto implica \SI{2.07}{kg/m^2} de muro; presupuestos FSPR de 2019 usan 1.03 a \SI{1.13}{kg/m^2}.
\item \textbf{Pendiente de cubierta.} La planta rotula 60\,\%; el corte, las elevaciones y el propio presupuesto (faldón de \SI{3.50}{m}) dan 40\,\%. La cubierta medida es \SI{\vTecho}{m^2}, sin traslapos según NCh 353.
\end{itemize}

\Needspace{382pt}% anticorte
\section{Otras inconsistencias del expediente}
\begin{itemize}[nosep]
\item Las EETT se contradicen en la cadena: 6\,$\phi$10 con estribos a 0.20 (punto 3.4) frente a $\phi$12 con estribos a 0.15 (punto 3.5). La planilla usa el punto 3.4; con 6\,$\phi$12 la masa por metro sube de 4.63 a \SI{6.57}{kg/m} (+42\,\%) y la partida 3.5 pasaría de 195 a unos \SI{276}{kg} (estimación por razón de masas, sin recalcular traslapos).
\item El sobrecimiento se especifica armado con 4\,$\phi$12 y estribos $\phi$6 a 0.20 (EETT 2.8), sin partida propia; puede estar incluido en el precio unitario, que es igual al de fundaciones armadas.
\item El ladrillo de las EETT mide 14\,cm de ancho y el muro se dibuja de 15\,cm. La ventana chica mide 0.80\,m de alto en elevación y la partida 15.2 dice 40\,$\times$\,60.
\item Radier, poliestireno y cielo usan \SI{53.3}{m^2}; el interior medido es \SI{51.03}{m^2} y la diferencia no se reproduce con ningún criterio encontrado.
\item Las EETT piden $\phi$10 en encuentros de muros; NCh1928 6.4.3.5 exige al menos una barra de \SI{12}{mm} en cada borde libre y encuentro. Se informa como consulta al proyectista, sin rediseñar. Las cuantías mínimas de 6.4.3.2 (0.06\,\% por dirección y 0.15\,\% total) se cumplen (0.09\,\% vertical y 0.07\,\% horizontal), y la escalerilla de \SI{4}{mm} cumple 6.1.2 (diámetro no mayor que la mitad de la junta de \SI{10}{mm}).
\item NCh1928 7.1 exige llenar con hormigón de relleno todo hueco que aloja armadura. Ni el presupuesto ni las EETT tienen esa partida; su volumen depende de los huecos del ladrillo, que el expediente no define, por lo que no se cubicó. NCh1928 6.4.3.6 pide además refuerzo horizontal en la base y el remate de los vanos, que tampoco aparece.
\item El acero especificado es A440-280H. Con hormigón G25, ACI 318-19 (25.4.2.3 y 25.5.2.1) da un largo de desarrollo de 27\,$d$ y un traslapo clase B de 35\,$d$, de modo que los 40\,$d$ de las EETT y de NCh1928 6.1.5 bastan. Con A630-420H se necesitarían 52\,$d$. NCh430 adopta ACI 318-05; que la ecuación coincida está [Por verificar].
\item Las EETT designan los hormigones como H-5 y H-25, designación anterior a NCh170:2016 (grados G por probeta cilíndrica); para cotizar hay que fijar el grado G con el proveedor.
\item Las EETT remiten a planos de estructuras y fundaciones que el expediente no incluye.
\end{itemize}

\section{Correcciones a la cubicación propia}
\begin{juntos}
La primera versión de esta cubicación tenía errores que la revisión en frío y la lectura de NCh 353 detectaron. Dos de ellos fabricaban coincidencias falsas con el presupuesto (relleno y cerámica), lo que muestra por qué una coincidencia no se acepta sin revisar la base.
\begin{table}
\centering\small
\begin{tabularx}{\linewidth}{lp{2.9cm}S[table-format=3.2]S[table-format=3.2]Xp{2.3cm}}
\toprule
Ítem & Partida & {Antes} & {Ahora} & Motivo & Detectado por\\
\midrule
2.4 & Relleno estabilizado e = 0.15 m & \num{9.36} & \num{7.65} & Relleno sobre la huella; EETT 2.4 lo define como base del radier (interior) & Revisor\\
11.1 & Cerámica 30 x 30 & \num{51.03} & \num{7.95} & Cerámica en todo el interior; EETT 11.2: solo cocina y baño & Revisor\\
3.3 & Tensores & \num{164.84} & \num{136.27} & Doble conteo de $\phi$8; faltaban +5\,\% de patas y la penetración de 40\,$d$ en fundación (NCh 353 6.2.8) & Revisor, ciega, NCh 353\\
3.2 & Escalerilla & \num{100.16} & \num{67.12} & Sin descontar vanos y con alambre de 4.2 mm (EETT: 4 mm) & Revisor\\
3.6 & Moldaje cadena & \num{22.48} & \num{23.77} & Faltaban fondos bajo dinteles y tapas de extremos & Revisor\\
8.1 & Bajadas AALL & \num{12.60} & \num{15.60} & Se tomó la cota del alero como largo de bajada & Revisor\\
3.1 & Ladrillo gran titán & \num{78.07} & \num{88.56} & Vanos descontados al 100\,\%; NCh 353 no descuenta los menores a \SI{1.5}{m^2} & NCh 353, revisor\\
3.5 & Enfierradura cadena & \num{173.65} & \num{194.77} & Sin traslapos de 40\,$d$ ni +5\,\% de patas & NCh 353, revisor\\
\bottomrule
\end{tabularx}
\caption{Correcciones a la cubicación propia tras la revisión en frío y la lectura de NCh 353.}
\label{tab:corr}
\end{table}

Todas las correcciones quedaron registradas en la hoja Conciliacion de la planilla, con el valor anterior y el nuevo.
\end{juntos}

\section{Supuestos y límites}
\label{sec:supuestos}
\begin{itemize}[nosep]
\item Ladrillo de NPT $\pm0.00$ a +2.40 y cadena de $15\times30$ de +2.40 a +2.70 (elevaciones y EETT 3.4). El corte dibuja el muro hasta cerca de +2.94; si el ladrillo llegara a +2.70, el ladrillo subiría a \SI{99.8}{m^2}.
\item Fundación $0.40\times0.60$, sobrecimiento $0.15\times0.20$ y emplantillado de \SI{5}{cm}, inferidos a falta de planos de estructuras.
\item Hilada de \SI{8.1}{cm} (ladrillo de 7.1 y junta de \SI{1.0}{cm}); escalerilla de dos alambres de \SI{4}{mm} con travesaños de \SI{8.5}{cm} cada \SI{25}{cm}, a confirmar con la ficha del fabricante.
\item Tensores: 38 de $\phi$8 y 40 de $\phi$10 contados panel por panel, con penetración de 40\,$d$ en la fundación (NCh353 6.2.8): \SI{3.22}{m} y \SI{3.30}{m}. Bajadas: 4 de \SI{3.90}{m}. Vanos sin pilar.
\item No se revisaron cerchas, terminaciones restantes ni precios unitarios.
\end{itemize}
Cubicación referencial elaborada desde los documentos públicos indicados, como ejercicio independiente. No reemplaza un proyecto ni un presupuesto firmado por un profesional competente. La planilla de respaldo (SedeC\_cubicacion\_v6.xlsx) contiene los datos con su fuente, las fórmulas, la conciliación y los controles.

\end{document}
